\section{Computational results}
\label{sec:comp}
% Every number to be filled in is written as \textbf{??}.  The comment
% "% [Name]" next to it is the name printed by  python make_tables.py ,
% which lists all values at the end of its output.

Since no other method is known for the IQCP with simultaneous modification of
lead times and capacities, we compare our algorithms with the general-purpose
approach of solving a mathematical programming model of the problem, and we
measure how their running times grow with the size of the tree.

\subsection{Setting}

Algorithms~\ref{alg1}, \ref{alg:A2}, \ref{alg3} and~\ref{alg4} were
implemented in Python~\textbf{??} % [RPython]
as stated, with $\varepsilon=10^{-6}$ in the augmentation setting; the
implementation reproduces Examples~\ref{exp1}--\ref{exp4} exactly. All runs
used a single core of an \textbf{??} % [RCpu]
with \textbf{??}\,GB of memory. % [RRam]
The code and instances are available at
\texttt{https://github.com/XXXX/iqcp-benchmark}. % repository URL

All instances use $\sigma=60$, lead times in $\{5,\dots,40\}$, capacities in
$\{1,\dots,15\}$, unit costs $\alpha^\pm\in\{5,\dots,40\}$ and
$\beta^\pm\in\{50,\dots,300\}$, $\bar x^-=\ell\,U(0.5,1)$,
$\bar y^-=c\,U(0.3,0.9)$, $\bar x^+=\ell\,U(0.5,2)$ and $\bar y^+=c\,U(1,4)$ on
$T_1$ ($\bar y^+=0$ on $T_2$); instances in which $v^*$ is already a center are
redrawn. We use random recursive trees (\emph{rrt}, paths of length
$\Theta(\log n)$) and caterpillars with two spines leaving $v^*$ (\emph{cat},
paths of length $\Theta(n)$). In the variants \emph{rrt-hc} and \emph{cat-hc}
all capacities of $T_2$ exceed those of $T_1$, which by
Lemma~\ref{lem:filter} is the configuration in which
Assumption~\ref{assumption2} is attainable. The family \emph{catA2} is the
worst case of Algorithm~\ref{alg:A2}: the spine capacities of $T_1$ decrease
strictly, so that $\mathcal B^c$ has $\Theta(n)$ elements in
$[\underline\tau,\overline\tau]$.

\subsection{Comparison with mathematical programming solvers}
\label{subsec:comp-gurobi}

We model the full IQCP without using any structural result of this paper:
all edges may be modified and any branch of $v^*$ may become the primary
subtree. With $w(e)=\sigma/\tilde c(e)$, the quickest distance of a leaf is
$\tilde\ell(P_{v^*t})+\max_{e\in P_{v^*t}}w(e)$; the maxima are modelled with
binary variables and Theorem~\ref{theo1} with binaries selecting the
assumption, the primary branch and $t^{**}$. Imposing
$\tilde c(e)\,w(e)=\sigma$ exactly gives a nonconvex mixed-integer
quadratically constrained programme (MIQCP), which we solve to global
optimality with Gurobi~\textbf{??}. % [RGurobi]
Since the bilinear constraints cannot be handled exactly by linear
mixed-integer solvers, Gurobi is the natural exact benchmark. It ran on one
thread with a time limit of \textbf{??}\,s, % [RTimeLimit]
on \textbf{??} instances per family and size, in both settings. % [RRepsSolv]

\begin{table}[htbp]
\centering
\small
\caption{Our algorithms against Gurobi on the exact MIQCP (time limit
\textbf{??}\,s, all families pooled). Times are means; \#opt: runs Gurobi
solved to proven optimality, and its time is the mean over these runs; ratio:
median of Gurobi's time divided by ours over these runs; gap: largest final
MIP gap of the runs that reached the time limit; $\Delta$: largest relative
deviation between Gurobi's optimum and ours.}
\label{tab:solvers}
% Copy the rows from VALUES_FOR_PAPER.txt (or latex/table_solvers.tex).
\begin{tabular}{lrrr|rrrrr}
\hline
 & & & Ours & \multicolumn{5}{c}{Gurobi (MIQCP)} \\
Setting & $n$ & runs & [ms] & \#opt & [s] & ratio & gap & $\Delta$ \\
\hline
Reduction    &   17 & & & & & & & \\
             &   33 & & & & & & & \\
             &   65 & & & & & & & \\
             &  129 & & & & & & & \\
             &  257 & & & & & & & \\
             &  513 & & & & & & & \\
             & 1025 & & & & & & & \\
             & 2049 & & & & & & & \\
\hline
Augmentation &   17 & & & & & & & \\
             &   33 & & & & & & & \\
             &   65 & & & & & & & \\
             &  129 & & & & & & & \\
             &  257 & & & & & & & \\
             &  513 & & & & & & & \\
             & 1025 & & & & & & & \\
             & 2049 & & & & & & & \\
\hline
\end{tabular}
\end{table}

\begin{figure}[htbp]
\centering
% Replace the box by  \includegraphics[width=\textwidth]{iqcp_scaling.pdf}
% after uploading latex/iqcp_scaling.pdf.
\fbox{\parbox[c][5cm][c]{0.9\textwidth}{\centering Figure: iqcp\_scaling.pdf}}
\caption{(a) Running time of the four algorithms (colour: algorithm, line
style: family; reference slopes $2$ and $3$). (b) Our algorithms and Gurobi on
the MIQCP (solid: reduction, dashed: augmentation; dotted: time limit).
(c) Evaluations of $F_i^{+}$ in Algorithm~\ref{alg3} against
$\log_{10}(1/\varepsilon)$, $n=801$.}
\label{fig:scaling}
\end{figure}

Table~\ref{tab:solvers} and Figure~\ref{fig:scaling}(b) summarise the
comparison. On all \textbf{??} runs that Gurobi solved to optimality, % [RgrbOptimal]
our value coincides with Gurobi's up to a relative deviation of
\textbf{??}, % [RgrbMaxDev]
including the \textbf{??} runs whose optimum is attained under % [RgrbAtwoWins]
Assumption~\ref{assumption2}. Gurobi solved all runs up to
$n=\textbf{??}$ (reduction) and $n=\textbf{??}$ (augmentation), % [RgrbAllSolvedRed], [RgrbAllSolvedAug]
where it was slower than our algorithms by median factors of
\textbf{??} and \textbf{??}, respectively; % [RspeedupRed], [RspeedupAug]
on larger trees it reached the time limit with final gaps of up to
\textbf{??}\,\%, % [RgrbMaxGap]
while our algorithms solved every instance exactly in a fraction of a second.

\subsection{Growth with the size of the instance}
\label{subsec:comp-scaling}

\begin{table}[htbp]
\centering
\small
\caption{Mean running time in milliseconds (\textbf{??} instances per entry) % [RReps]
and exponent $p$ of a least-squares fit $t\approx Cn^{p}$. For \emph{catA2}
the sizes are $41,81,161,321,641,1281$.}
\label{tab:runtimes}
% Copy the rows from latex/table_runtimes.tex, or fill in by hand.
\begin{tabular}{ll rrrrrrr r}
\hline
Algorithm & Family & \multicolumn{7}{c}{$n$} & $p$ \\
& & 101 & 201 & 401 & 801 & 1601 & 3201 & 6401 & \\
\hline
Alg.~\ref{alg1}  & \emph{rrt}    & & & & & & & & \\
Alg.~\ref{alg1}  & \emph{cat}    & & & & & & & & \\
Alg.~\ref{alg:A2} & \emph{rrt-hc} & & & & & & & & \\
Alg.~\ref{alg:A2} & \emph{cat-hc} & & & & & & & & \\
Alg.~\ref{alg:A2} & \emph{catA2}  & & & & & & & & \\
Alg.~\ref{alg3}  & \emph{rrt}    & & & & & & & & \\
Alg.~\ref{alg3}  & \emph{cat}    & & & & & & & & \\
Alg.~\ref{alg4}  & \emph{rrt-hc} & & & & & & & & \\
Alg.~\ref{alg4}  & \emph{cat-hc} & & & & & & & & \\
\hline
\end{tabular}
\end{table}

Table~\ref{tab:runtimes} and Figure~\ref{fig:scaling}(a) report running times
up to $n=6401$ with the exponent $p$ of a fit $t\approx Cn^{p}$. Apart from
\emph{catA2}, the exponents lie between \textbf{??} and \textbf{??}, % [RexpMin], [RexpMax]
below the worst-case bounds of Sections~\ref{sec3} and~\ref{sec4}, because on
random trees the dynamic programmes carry far fewer than $n$ breakpoints and
few leaves pass the screening of \eqref{eq:prefix}. On \emph{catA2} the
number of evaluations of \eqref{eq:D1} grows linearly and the exponent is
\textbf{??} against the bound $O(n^3)$; % [RexpAtwoCatAtwo]
the number $K$ of additional breakpoints of $D_1^-$ never exceeded
\textbf{??}, % [RKmax]
so the condition $K=O(n)$ of Theorem~\ref{theo:A2minus} held on all
instances. Finally, the number of evaluations in the ternary search grows
linearly in $\log(1/\varepsilon)$ (Figure~\ref{fig:scaling}(c)), and lowering
$\varepsilon$ from $10^{-2}$ to $10^{-8}$ increases the running time only by a
factor of \textbf{??}. % [RepsTimeRatio]
